\section{The Hamiltonian}

\SourcePage{41}{44}
In quantum mechanics the wave function $\Psi$ fully specifies the state of a
physical system. This means that prescribing $\Psi$ at some instant not
only describes all the properties of the system at that instant, but also
determines its behavior at every later instant---of course, only to the degree
of completeness that quantum mechanics permits at all. Mathematically, this
means that at each instant the value of the time derivative
$\partial\Psi/\partial t$ must be determined by the value of $\Psi$ itself at
that same instant, and, by the superposition principle, this dependence must be
linear. In the most general form we may write
\begin{equation}
  i\hbar\frac{\partial\Psi}{\partial t}=\hat H\Psi,
  \tag{8.1}
\end{equation}
where $\hat H$ is a certain linear operator; the factor $i\hbar$ has been
introduced here for a reason that will become clear below.

Since the integral $\int\Psi^*\Psi\,dq$ is constant in time, we have
\[
  \frac{d}{dt}\int|\Psi|^2\,dq
  =\int\frac{\partial\Psi^*}{\partial t}\Psi\,dq
  +\int\Psi^*\frac{\partial\Psi}{\partial t}\,dq=0.
\]
Substituting (8.1) here and applying the definition of the transposed operator
in the first integral, we obtain (after suppressing the common factor
$i/\hbar$)
\[
  \begin{aligned}
    \int\Psi\hat H^*\Psi^*\,dq-\int\Psi^*\hat H\Psi\,dq
      &={}\int\Psi^*\tilde H^*\Psi\,dq
        -\int\Psi^*\hat H\Psi\,dq\\
      &=\int\Psi^*(\tilde H^*-\hat H)\Psi\,dq=0.
  \end{aligned}
\]
Because this equality must hold for an arbitrary function $\Psi$, it follows
that identically $\hat H^+=\hat H$; the operator $\hat H$ is therefore
Hermitian.
\SourcePage{42}{45}
Let us determine the physical quantity to which it corresponds. We use the
limiting expression (6.1) for the wave function and write
\[
  \frac{\partial\Psi}{\partial t}
  =\frac{i}{\hbar}\frac{\partial S}{\partial t}\Psi
\]
(the slowly varying amplitude $a$ need not be differentiated). Comparing this
equality with definition (8.1), we see that in the limiting case $\hat H$
reduces to multiplication by $-\partial S/\partial t$. Thus this last quantity
is the physical quantity into which the Hermitian operator $\hat H$ passes.

But $-\partial S/\partial t$ is precisely the Hamilton function $H$ of the
mechanical system. Hence $\hat H$ is the operator corresponding in quantum
mechanics to the Hamilton function. This operator is termed the
\textit{Hamiltonian operator}, or simply the \textit{Hamiltonian} of the
system. Once the form of the Hamiltonian is known, equation (8.1) fixes the wave functions of the given
physical system. This fundamental equation of quantum mechanics is called the
\textit{wave equation}.
